\documentclass[11pt,a4paper]{article}
\usepackage[utf8]{inputenc}
\usepackage[spanish]{babel}
\usepackage{amsmath,amsfonts,amssymb}
\usepackage{geometry}
\geometry{top=2.5cm,bottom=2.5cm,left=2.5cm,right=2.5cm}
\usepackage{hyperref}
\usepackage{booktabs}
\usepackage{xcolor}
\usepackage{tcolorbox}
\usepackage{enumitem}

% Definición de colores institucionales
\definecolor{externadoDarkBlue}{RGB}{10, 34, 64}
\definecolor{externadoAccent}{RGB}{0, 114, 206}
\definecolor{lightGray}{RGB}{245, 247, 250}

\hypersetup{
    colorlinks=true,
    linkcolor=externadoAccent,
    urlcolor=externadoAccent,
    titlecolor=externadoDarkBlue
}

\tcbset{
    colback=lightGray,
    colframe=externadoDarkBlue,
    fonttitle=\bfseries\color{white},
    arc=3mm
}

\title{
    \vspace{-1cm}
    \textbf{\Large UNIVERSIDAD EXTERNADO DE COLOMBIA}\\
    \textbf{\large PREGRADO EN CIENCIA DE DATOS --- EMPRENDIMIENTO}\\
    \vspace{0.5cm}
    \Large \textbf{TALLER 2: CONTRATO DE EQUIPO (TERM SHEET)}\\
    \large \textbf{Proyecto: StatCredit AI --- Scoring Crediticio Alternativo B2B}
}
\author{\textbf{Grupo 5:} Santiago Sandoval | Sebastián Ramos | Julián Duarte | Tomás Rincón | Julián Jiménez}
\date{Semestre 2026-II}

\begin{document}

\maketitle

\hrule
\vspace{0.5cm}

\section{Objetivo del Taller}
El presente taller tiene como objetivo co-construir los lineamientos fundamentales del contrato de aprendizaje, el contrato de equipo y la declaración de expectativas y compromisos para el desarrollo del proyecto \textbf{StatCredit AI}. A través de este \textit{Term Sheet}, el equipo establece un pacto de responsabilidad ética, validación cuantitativa rigurosa y colaboración estratégica sostenida durante el semestre.

\section{Justificación e Importancia del Trabajo en Equipo}
El trabajo interdisciplinario en el pregrado de Ciencia de Datos constituye una competencia crítica para resolver problemas complejos en el sector financiero. La co-creación estratégica permite incorporar múltiples perspectivas (analítica de datos, ingeniería financiera y MLOps) para potenciar la innovación disruptiva (\textit{Tucker \& Abbasi, 2012}). Asimismo, la eficacia del desempeño operativo es directamente proporcional a la confianza, transparencia y empatía subyacentes en el equipo (\textit{Katzenbach \& Smith, 1993}).

De acuerdo con el informe de inversión de KPMG (2023), los comités de evaluación y fondos de venture capital priorizan equipos de gestión cohesionados, con habilidades técnicas demostradas y planes de ejecución claros por encima de ideas aisladas.

\section{Matriz de Integrantes e Ikigai del Equipo (Grupo 5)}

\begin{table}[h!]
\centering
\small
\begin{tabular}{p{3.5cm} p{4cm} p{4cm} p{3.5cm}}
\toprule
\textbf{Integrante} & \textbf{Fortaleza Principal (Ikigai)} & \textbf{Rol en StatCredit AI} & \textbf{Compromiso Clave} \\
\midrule
\textbf{Santiago Sandoval} & Modelado predictivo, MLOps e inferencia causal bayesiana. & Científico de Datos Lead / ML Architect & Garantizar auditabilidad y métricas de precisión algorítmica. \\
\midrule
\textbf{Sebastián Ramos} & Estrategia B2B, finanzas cuantitativas y modelos de negocio. & Director de Estrategia \& Negocios B2B & Definir modelo de monetización SaaS y proyecciones TAM/SAM/SOM. \\
\midrule
\textbf{Julián Duarte} & Arquitectura de software, APIs REST y gestión de datos. & Data Engineer / Backend Architect & Diseñar la tubería de ingesta y despliegue de la API. \\
\midrule
\textbf{Tomás Rincón} & Análisis de riesgo crediticio, regulación SARC y BI. & Analista de Riesgo \& Compliance & Asegurar cumplimiento normativo SARC y explicabilidad SHAP. \\
\midrule
\textbf{Julián Jiménez} & Desarrollo de software, interfaces y pruebas unitarias. & Software Engineer / QA & Garantizar la integridad operativa y pruebas de integración. \\
\bottomrule
\end{tabular}
\caption{Matriz de roles y complementariedad del equipo Grupo 5.}
\end{table}

\section{Pacto de Compromiso Ético y Protocolo de Manejo de Conflictos}

\begin{tcolorbox}[title=Acuerdos Fundamentales de Equipo (Term Sheet)]
\begin{enumerate}[leftmargin=*]
    \textbf{1. Rigor Técnico y Datos Reales:} Ninguna métrica predictiva o modelo de scoring será presentado sin validación objetiva. Se prohíbe el uso de datos sintéticos no justificados o la simulación de resultados sin sustento metodológico.
    \textbf{2. Cumplimiento de Entregas y Puntualidad:} Todos los avances serán consolidados en formato \LaTeX{} al menos 24 horas antes de la fecha límite estipulada. Las sesiones de trabajo requerirán puntualidad estricta.
    \textbf{3. Protocolo de Resolución de Desacuerdos:} Frente a divergencias técnicas o de visión estratégica, se aplicará el siguiente procedimiento paso a paso:
    \begin{itemize}
        \item \textit{Fase 1 (Argumentación Causal):} Cada postura debe sustentarse en datos objetivos, literatura académica o exigencias regulatorias (SARC).
        \item \textit{Fase 2 (Votación Ponderada):} Si persiste la paridad, el líder del área técnica correspondiente (Santiago en Ciencia de Datos, Sebastián en Estrategia) tendrá voto dirimente.
        \item \textit{Fase 3 (Consulta Docente):} Escalación al profesor de la asignatura únicamente tras agotar las instancias internas.
    \end{itemize}
    \textbf{4. Transparencia y Coevaluación:} La coevaluación del equipo reflejará fielmente la contribución individual, el cumplimiento de tareas y el respeto por los acuerdos del presente contrato.
\end{enumerate}
\end{tcolorbox}

\section{Firma e Identificación del Equipo}

\vspace{1cm}
\begin{table}[h!]
\centering
\begin{tabular}{ccc}
\underline{\hspace{4.5cm}} & \hspace{0.5cm} & \underline{\hspace{4.5cm}} \\
Santiago Sandoval & & Sebastián Ramos \\
Científico de Datos Lead & & Director de Estrategia B2B \\
\vspace{1cm} & & \\
\underline{\hspace{4.5cm}} & \hspace{0.5cm} & \underline{\hspace{4.5cm}} \\
Julián Duarte & & Tomás Rincón \\
Data Engineer & & Analista de Riesgo SARC \\
\vspace{1cm} & & \\
\multicolumn{3}{c}{\underline{\hspace{4.5cm}}} \\
\multicolumn{3}{c}{Julián Jiménez} \\
\multicolumn{3}{c}{Software Engineer / QA}
\end{tabular}
\end{table}

\end{document}
